
\begin{exercise}[subtitle={Stochastic convergence sometimes implies \(L^p\) \Coffeecup}]
    \label{ex: p to Lp when bounded}
    Assume that \(X_n\) and \(X\) are bounded random variables. That is 
    \(\sup_n\|X_n\| \le c\) and \(\|X\|\le c\) with probability one for some \(c
    \in \real\). Show that then \(X_n \overset{p}\to X\) implies \(X_n
    \overset{L^p}\to X\).
    \begin{remark}
        This is a weak form of the dominated convergence theorem, where the constant
        is the dominating function.
    \end{remark}
\end{exercise}
\begin{solution}
    We have for every \(\epsilon > 0\)
    \[\begin{aligned}
        \E[\big]{\norm{X_n - X}^p}
        &= \E[\big]{\norm{X_n - X}^p \ind{\norm{X_n - X} < \epsilon}}
        + \E[\big]{\norm{X_n - X}^p \ind{\norm{X_n - X} \ge \epsilon}}
        \\
        &\le \epsilon^p + (2c)^p \prob[\big]{\norm{X_n - X} \ge \epsilon}.
    \end{aligned}
    \]
    Due to \(X_n \overset{p}\to X\) the second term converges to \(0\)
    and we thus have
    \[
        \limsup_{n\to \infty} \E[\big]{\norm{X_n - X}^p} \le \epsilon^p.
    \]
    Since this is true for any \(\epsilon > 0\) we have the claim.
\end{solution}
